\clearpage

% 3. METODO
\section{Método de Trabalho}

O Clube de Matemática combina aprendizado estruturado com prática colaborativa. Para
equilibrar teoria e aplicação, dois tipos complementares de sessões serão organizados:
\textbf{sessões de estudo} e \textbf{sessões de resolução de problemas}. Além das reuniões
regulares, o clube manterá um ambiente técnico compartilhado para comunicação, documentação
e troca de conhecimento.

\subsection{Sessões de Estudo}

\begin{infobox}
  \textbf{Frequência:} uma vez por semana \quad\textbf{Duração:} 2 horas
\end{infobox}

\subsubsection{Programa da Sessão}

\begin{itemize}
  \item \textbf{00:00 -- 00:10} | \textit{Abertura:} revisão do backlog e votação do que
    será discutido na sessão.
  \item \textbf{00:10 -- 01:40} | \textit{Exploração:} discussão e resolução colaborativa
    dos temas escolhidos. Quem trouxe um problema ou referência naturalmente contribui mais
    --- mas qualquer membro pode assumir a explicação.
  \item \textbf{01:40 -- 01:55} | \textit{Registro:} soluções e descobertas relevantes são
    documentadas no repositório.
  \item \textbf{01:55 -- 02:00} | \textit{Próxima sessão:} definição do facilitador
    rotativo para o próximo encontro.
\end{itemize}

\subsection{Facilitador Rotativo}

Para garantir que as sessões comecem e fluam bem, um membro assume o papel de
\textbf{facilitador} a cada semana. Esse papel é leve: garantir que a sessão comece no
horário, conduzir a votação do backlog e manter o tempo. O facilitador \textbf{não precisa
preparar conteúdo} --- o conteúdo vem do grupo.

A rotação é definida no final de cada sessão. Se o facilitador designado não puder
comparecer, qualquer outro membro assume voluntariamente.

\subsection{Ferramentas e Ambiente}

\begin{itemize}
  \item Um canal de comunicação no Slack (ou plataforma equivalente da 42) para coordenação
    e discussões;
  \item Um repositório compartilhado (Vogsphere ou GitHub) para soluções, anotações e
    materiais de referência.
\end{itemize}

\subsection{Registro de Soluções e Compartilhamento de Conhecimento}

Documentar soluções de problemas e descobertas matemáticas é parte importante do processo
de aprendizado. Os membros são incentivados a criar registros após resolver problemas
desafiadores ou descobrir uma técnica nova.
